\documentclass[10pt]{beamer}
\usetheme{metropolis}
\usepackage{graphicx}
\usepackage{listings}
\usepackage{booktabs}
\usepackage{xcolor}
\definecolor{codebg}{HTML}{F5F5F5}
\definecolor{codegreen}{HTML}{2E7D32}
\definecolor{codegray}{HTML}{757575}
\definecolor{codeblue}{HTML}{1565C0}
\lstdefinestyle{rstyle}{language=R, backgroundcolor=\color{codebg}, basicstyle=\ttfamily\scriptsize,
  keywordstyle=\color{codeblue}\bfseries, stringstyle=\color{codegreen},
  commentstyle=\color{codegray}\itshape, breaklines=true, frame=single,
  rulecolor=\color{codegray!30}, showstringspaces=false, tabsize=2,
  morekeywords={library,ggplot,aes,geom_boxplot,geom_violin,geom_jitter,geom_beeswarm,
    geom_point,stat_summary,scale_fill_manual,scale_color_manual,coord_flip,
    theme_minimal,theme,element_text,labs,position_jitter,outlier.shape}}
\lstdefinestyle{pythonstyle}{language=Python, backgroundcolor=\color{codebg}, basicstyle=\ttfamily\scriptsize,
  keywordstyle=\color{codeblue}\bfseries, stringstyle=\color{codegreen},
  commentstyle=\color{codegray}\itshape, breaklines=true, frame=single,
  rulecolor=\color{codegray!30}, showstringspaces=false, tabsize=4,
  morekeywords={import,from,as,pd,plt,sns,np,fig,ax,boxplot,violinplot,stripplot,swarmplot}}
\title{Distributions: Boxplots, Violins \& Beeswarm}
\subtitle{Workshop 3 --- Module 2}
\author{Prof.Asoc.\ Endri Raco}
\institute{Data Visualization for Data Scientists \\ UNYT Tirana}
\date{2025/26}
\begin{document}

\begin{frame}
  \titlepage
\end{frame}

\begin{frame}{Module Roadmap}
  \begin{enumerate}
    \item Boxplot anatomy: the five-number summary
    \item What the boxplot hides: shape blindness
    \item Violin plots: full density shape
    \item Strip, jitter, and beeswarm: showing raw data
    \item Raincloud plots: the best of all three
    \item Distribution chart selection guide
  \end{enumerate}
  \begin{exampleblock}{Learning Outcome}
    You will choose between boxplot, violin, strip, beeswarm, and
    raincloud based on sample size, audience, and analytical goal.
  \end{exampleblock}
\end{frame}

\section{Boxplot Anatomy}

\begin{frame}{The Five-Number Summary + IQR + Outliers}
  \begin{center}
    \includegraphics[width=0.78\textwidth]{figures_w03m02/boxplot_anatomy}
  \end{center}
\end{frame}

\begin{frame}{What the Boxplot Hides}
  \begin{center}
    \includegraphics[width=0.92\textwidth]{figures_w03m02/box_vs_hist}
  \end{center}
  \begin{alertblock}{Pro-Tip}
    Three very different distributions --- unimodal, bimodal, and
    right-skewed --- can produce nearly identical boxplots. The box
    compresses all shape information into five numbers. Always pair a
    boxplot with jitter, violin, or histogram to show the full picture.
  \end{alertblock}
\end{frame}

\begin{frame}[fragile]{Grouped Boxplot in R vs Python}
  \begin{columns}[T]
    \begin{column}{0.48\textwidth}
      \textbf{R (ggplot2)}
\begin{lstlisting}[style=rstyle]
ggplot(apps, aes(x = Category,
    y = Rating, fill = Category)) +
  geom_boxplot(
    width = 0.5,
    notch = TRUE,
    outlier.alpha = 0.2) +
  stat_summary(fun = mean,
    geom = "point",
    shape = 18, size = 3,
    color = "#E53935") +
  coord_flip() +
  theme_minimal() +
  theme(legend.position = "none")
\end{lstlisting}
    \end{column}
    \begin{column}{0.48\textwidth}
      \textbf{Python (seaborn)}
\begin{lstlisting}[style=pythonstyle]
sns.boxplot(data=apps,
    x="Category", y="Rating",
    palette="Set2",
    width=0.5, notch=True,
    flierprops={"alpha": 0.2})

# Or matplotlib:
bp = ax.boxplot(data_list,
    labels=cats,
    patch_artist=True,
    notch=True, widths=0.5)
# Add means:
ax.scatter(range(1,n+1), means,
    c="#E53935", s=40,
    marker="D", zorder=5)
\end{lstlisting}
    \end{column}
  \end{columns}
\end{frame}

\begin{frame}{Grouped Boxplot with Notch and Means}
  \begin{center}
    \includegraphics[width=0.82\textwidth]{figures_w03m02/grouped_boxplot}
  \end{center}
  \begin{exampleblock}{Data Insight}
    \textbf{Notches} show a confidence interval on the median: if two
    boxes' notches do not overlap, the medians are significantly different
    at approximately the 95\% level. \textbf{Diamond markers} show group
    means for comparison with medians.
  \end{exampleblock}
\end{frame}

\section{Violin Plots}

\begin{frame}{Boxplot vs Violin: What You Gain}
  \begin{center}
    \includegraphics[width=0.92\textwidth]{figures_w03m02/violin}
  \end{center}
\end{frame}

\begin{frame}[fragile]{Violin Plots in R vs Python}
  \begin{columns}[T]
    \begin{column}{0.48\textwidth}
      \textbf{R}
\begin{lstlisting}[style=rstyle]
ggplot(apps, aes(x = Type,
    y = Rating, fill = Type)) +
  geom_violin(
    alpha = 0.3,
    draw_quantiles = c(0.25,
      0.5, 0.75)) +
  geom_boxplot(
    width = 0.1,
    outlier.shape = NA,
    fill = "white") +
  scale_fill_manual(values =
    c(Free="#1565C0",
      Paid="#E53935")) +
  theme_minimal() +
  theme(legend.position = "none")
\end{lstlisting}
      {\small Embed a narrow boxplot inside the
      violin for both shape and summary.}
    \end{column}
    \begin{column}{0.48\textwidth}
      \textbf{Python}
\begin{lstlisting}[style=pythonstyle]
# seaborn (recommended)
sns.violinplot(data=apps,
    x="Type", y="Rating",
    hue="Type",
    palette={"Free":"#1565C0",
             "Paid":"#E53935"},
    inner="quartile",  # or "box"
    split=False, alpha=0.3)

# split=True for side-by-side:
sns.violinplot(data=apps,
    x="Category", y="Rating",
    hue="Type", split=True,
    inner="quart",
    palette={"Free":"#1565C0",
             "Paid":"#E53935"})
\end{lstlisting}
    \end{column}
  \end{columns}
\end{frame}

\begin{frame}{Split Violin: Two Groups Side by Side}
  \begin{center}
    \includegraphics[width=0.78\textwidth]{figures_w03m02/split_violin}
  \end{center}
\end{frame}

\section{Strip, Jitter \& Beeswarm}

\begin{frame}{Showing Individual Observations}
  \begin{center}
    \includegraphics[width=0.92\textwidth]{figures_w03m02/beeswarm}
  \end{center}
  \begin{exampleblock}{Data Insight}
    \textbf{Strip/jitter} adds random horizontal noise --- fast but
    imprecise. \textbf{Beeswarm} stacks points systematically to reveal
    local density without randomness. Both show every data point,
    which boxplots and violins hide.
  \end{exampleblock}
\end{frame}

\begin{frame}[fragile]{Jitter + Beeswarm in Code}
  \begin{columns}[T]
    \begin{column}{0.48\textwidth}
      \textbf{R}
\begin{lstlisting}[style=rstyle]
# Jitter
ggplot(apps, aes(x = Type,
    y = Rating)) +
  geom_boxplot(
    outlier.shape = NA,
    fill = "#E3F2FD") +
  geom_jitter(width = 0.15,
    height = 0, alpha = 0.2,
    size = 0.8,
    color = "#E53935") +
  theme_minimal()

# Beeswarm
# install.packages("ggbeeswarm")
library(ggbeeswarm)
ggplot(apps, aes(x = Type,
    y = Rating)) +
  geom_beeswarm(
    cex = 0.5, alpha = 0.3,
    color = "#1565C0", size = 0.8)
\end{lstlisting}
    \end{column}
    \begin{column}{0.48\textwidth}
      \textbf{Python}
\begin{lstlisting}[style=pythonstyle]
# Jitter (strip)
sns.stripplot(data=apps,
    x="Type", y="Rating",
    jitter=0.2, alpha=0.2,
    size=2, color="#E53935")

# Combined: box + strip
fig, ax = plt.subplots()
sns.boxplot(data=apps,
    x="Type", y="Rating",
    width=0.4, ax=ax,
    fliersize=0)
sns.stripplot(data=apps,
    x="Type", y="Rating",
    jitter=0.15, alpha=0.2,
    size=2, color="#E53935",
    ax=ax)

# Beeswarm (seaborn)
sns.swarmplot(data=apps,
    x="Type", y="Rating",
    size=2, alpha=0.3)
\end{lstlisting}
    \end{column}
  \end{columns}
\end{frame}

\section{Raincloud Plots}

\begin{frame}{Raincloud: Box + Violin + Jitter Combined}
  \begin{center}
    \includegraphics[width=0.78\textwidth]{figures_w03m02/raincloud}
  \end{center}
  \begin{alertblock}{Pro-Tip}
    The raincloud plot (Allen et al., 2019) is the gold standard for
    distribution visualisation in scientific publications. It shows
    summary statistics (box), density shape (half-violin), and raw data
    (jitter) simultaneously. In R: \texttt{ggrain::geom\_rain()}. In
    Python: manual composition with seaborn.
  \end{alertblock}
\end{frame}

\section{Chart Selection}

\begin{frame}{Distribution Chart Comparison}
  \begin{center}
    \includegraphics[width=0.92\textwidth]{figures_w03m02/comparison_card}
  \end{center}
\end{frame}

\begin{frame}{When to Use Which}
  \centering
  \begin{tabular}{lll}
    \toprule
    \textbf{Scenario} & \textbf{Best Chart} & \textbf{Why} \\
    \midrule
    Many groups ($>$8) & Boxplot & Compact summary \\
    Shape matters & Violin & Shows bimodality \\
    Small n ($<$200) & Strip / Beeswarm & Every point visible \\
    Publication & Raincloud & Maximum information \\
    Quick EDA & Boxplot + jitter & Fast, informative \\
    Two-group comparison & Split violin & Side-by-side shape \\
    \bottomrule
  \end{tabular}
\end{frame}

\section{Complete Examples}

\begin{frame}[fragile]{R: Boxplot + Jitter (Publication Ready)}
  \begin{columns}[T]
    \begin{column}{0.52\textwidth}
\begin{lstlisting}[style=rstyle]
ggplot(apps |>
    filter(Category %in% top4),
  aes(x = fct_reorder(Category,
    Rating, .fun = median),
    y = Rating, fill = Category)) +
  geom_boxplot(
    width = 0.4,
    outlier.shape = NA,
    alpha = 0.3) +
  geom_jitter(
    width = 0.1, alpha = 0.15,
    size = 0.5,
    color = "#E53935") +
  stat_summary(fun = mean,
    geom = "point",
    shape = 18, size = 3) +
  coord_flip() +
  theme_minimal() +
  theme(legend.position = "none") +
  labs(title = "Rating by Category",
       x = NULL, y = "Rating")
\end{lstlisting}
    \end{column}
    \begin{column}{0.45\textwidth}
      \includegraphics[width=\textwidth]{figures_w03m02/r_output}
    \end{column}
  \end{columns}
\end{frame}

\begin{frame}[fragile]{Python: Violin + Strip}
  \begin{columns}[T]
    \begin{column}{0.52\textwidth}
\begin{lstlisting}[style=pythonstyle]
fig, ax = plt.subplots(figsize=(6,5))

# Violin layer
parts = ax.violinplot(
    data_list,
    positions=range(1, 5),
    showmeans=True,
    showmedians=True,
    widths=0.6)
for pc, c in zip(
    parts["bodies"], colors):
    pc.set_facecolor(c)
    pc.set_alpha(0.25)

# Strip layer (jitter overlay)
for i, d in enumerate(data_list):
    xj = np.random.normal(
        i + 1, 0.04, min(60, len(d)))
    ax.scatter(xj,
        np.random.choice(d, 60),
        s=4, c=colors[i],
        alpha=0.4, edgecolors="none")

ax.set_title("Violin + Strip")
\end{lstlisting}
    \end{column}
    \begin{column}{0.45\textwidth}
      \includegraphics[width=\textwidth]{figures_w03m02/py_output}
    \end{column}
  \end{columns}
\end{frame}

\section{Synthesis}

\begin{frame}{Key Takeaways}
  \begin{enumerate}
    \item Boxplots show the \textbf{five-number summary + outliers} but
          hide distribution shape (bimodality, skewness).
    \item \textbf{Violin plots} reveal the full density shape; embed a
          narrow boxplot inside for both shape and summary.
    \item \textbf{Strip/jitter} shows every data point; use for small
          $n < 200$. \textbf{Beeswarm} avoids random overlap.
    \item \textbf{Raincloud} (half-violin + box + jitter) is the publication
          gold standard (Allen et al., 2019).
    \item Always \textbf{pair boxplots with jitter or violin} to avoid
          the shape-blindness problem.
    \item In R: \texttt{geom\_boxplot + geom\_jitter}, \texttt{geom\_violin},
          \texttt{ggbeeswarm}. In Python: \texttt{sns.boxplot + stripplot},
          \texttt{sns.violinplot}, \texttt{sns.swarmplot}.
  \end{enumerate}
\end{frame}

\begin{frame}{Essential Readings}
  \begin{itemize}
    \item Allen, M. et al. (2019). ``Raincloud plots: a multi-platform tool
          for robust data visualization.'' \emph{Wellcome Open Research}, 4, 63.
    \item Wickham, H. \& Stryjewski, L. (2011). ``40 years of boxplots.''
          \emph{The American Statistician}.
    \item Wilke, C.~O. (2019). \emph{Fundamentals of Data Visualization},
          Chapters 7, 9 (Distributions).
    \item Hintze, J.~L. \& Nelson, R.~D. (1998). ``Violin plots: a box
          plot-density trace synergism.'' \emph{The American Statistician}.
  \end{itemize}
  \vspace{6pt}
  \textbf{Next Module}: Comparisons --- Bar Charts, Lollipop \& Cleveland Dot (W03-M03)
\end{frame}

\end{document}
